% =====================================================================
%  R. Nielsen, THE RIGHT OF THE STRONGER.  Speech at the Reformation Festival 1863 for the University.
%  Copenhagen: Gyldendal (F. Hegel), 1863.  15 pp.  (Den Stærkeres Ret. Tale ved Reformationsfesten 1863 for
%  Universitetet.)
%  ENGLISH TRANSLATION, from transcription.tex (Danish), 2026-10-08, made in the main session from the collated and
%  page-checked transcription.  Conventions: TRANSLATION-PLAYBOOK.md; terms as in the earlier Nielsen translations
%  (ebook-method/skjaebne-og-forsyn/GLOSSARY.md).  Page numbers of the 1863 print in the margin.
%  Terms: Ret = right; Magt = power; Uret = wrong; Afmagt = impotence; Retfærdighed = justice; Dom = judgment;
%  Verdensdom = world-judgment.  The refrain „den Stærkere vil have Ret“ = "the stronger will have right".
%  Scripture and Luther: Nielsen's Danish, in familiar English wording where his Danish matches it (Prov. 16:32;
%  Matt. 19:16; Luther at Worms, "Here I stand, I can do no other"); „vor faste Borg“ echoes Luther's hymn
%  „Ein' feste Burg“ ("A mighty fortress").  Misprints (marked PRINTED AS IS in the Danish) are translated as meant.
% =====================================================================
\documentclass[12pt,a4paper,openany]{book}
\usepackage[utf8]{inputenc}
\usepackage[T1]{fontenc}
\usepackage{libertinus}
\usepackage{libertinust1math}
\usepackage{textalpha}
\usepackage{graphicx}
\usepackage[english]{babel}
\usepackage{enumitem}
\usepackage[protrusion=true,expansion=true]{microtype}
\usepackage{setspace}
\usepackage[margin=1.3in]{geometry}
\usepackage{fancyhdr}
\usepackage{hyperref}
\hypersetup{pdftitle={The Right of the Stronger}, pdfauthor={Rasmus Nielsen}, hidelinks}
\setlength{\headheight}{15pt}
\pagestyle{fancy}
\fancyhf{}
\fancyhead[LE,RO]{\thepage}
\fancyhead[C]{\textit{R. Nielsen: The Right of the Stronger (1863)}}
\setstretch{1.3}
\usepackage{marginnote}
\newcommand{\opage}[1]{\marginnote{\footnotesize\textit{[#1]}}}

\begin{document}

% PDF 7, the title page, p. [1]
\begin{titlepage}
\begin{center}
\vspace*{4em}
{\Huge\textbf{The Right of the Stronger.}}\\[2em]
{\Large\textls[200]{Speech}}\\[1em]
{\large at the Reformation Festival 1863}\\[0.8em]
{\small for}\\[0.8em]
{\large\textls[200]{the University}}\\[0.8em]
{\small by}\\[0.8em]
{\large\textbf{R. Nielsen.}}\\[8em]
{\small\textbf{Copenhagen.}}\\[0.5em]
{\small Published by the Gyldendal Bookshop (F. Hegel).}\\[0.5em]
{\footnotesize\textls[200]{Thiele's Printing House.}}\\[0.5em]
\textit{1863.}\\[4em]
{\small Translated from the Danish}
\end{center}
\end{titlepage}

% ===== TEXT (pp. 3--15) =====
% --- p. 3 ---
\opage{3}On the way of the Reformation, in its strife, its course,
its struggle and victory, one thought shines out --- the old thought of
strength and weakness, of right and power; and this thought is
a judgment, a world-judgment on right and power: the stronger
will have right, the weaker must give way.

Yes, it is a world-judgment: the stronger will have right.
``We will have right, for we have power!'' So it sounded in threats
from Rome; ``we shall surely get power, as surely as we have right!''
So it sounded in reply from Wittenberg. Thus from every
side are heard the combatants' cries for right, for power, and for
a settlement between right and power; it is as if they all wanted the
same thing. And yet here is a difference, an essential difference.
For it does make an essential difference whether we understand
the unity of right with power so that it is the sensuous,
tangible power that, without further ado, is right, or
so that it is the holy, the incorruptible right that in
its essence is eternally power; it makes a world of difference whether
we judge strength so that the stronger is in the right
merely because he is strong, or so that he is in truth
the stronger only in so far as he has right.

Here there is a difference in right and here there is a difference in power:
here is physical power, and here is spiritual power; but whichever
of these powers one would turn into right, if that does
not mean that one would at the same time turn right into
power, then power is abused all the same.

% --- p. 4 ---
\opage{4}There is misery in the world, there is wretchedness and want,
when physical power is abused; for it is hard,
indescribably hard, to bow before brute power; it is hard,
indescribably hard, to acknowledge the right of the stronger when the
stronger is only a robber; it is hard, indescribably hard,
to bear the yoke of oppression, and hard to weigh freedom
against gold. And freedom weighed against gold! what then is the weight
of gold, when a Brennus throws his sword into the scale?
\textit{væ victis!} --- yes, woe to the vanquished.

And yet the misery is greater still when spiritual
power is abused; when deceit overshadows the
holy, when souls are caught and thoughts are bound and truth
fades and gives a light like the sun's dim light in an
eclipse. War with its horrors is indeed dreadful;
but the plague is more dreadful still; and slavery
in outward things is indeed a sad thing, but slavery in inward things,
spiritual slavery, the slavery of conscience and of thought,
is the plague itself, the sinister plague: this black death is perdition.

We may not now say that the Catholic Church began by
bending right and abusing power; for, temporally seen, it did not begin
with power. It began not in pomp
and glory, but in renunciation, in self-denial and outward
lowliness. It did not begin with a vicar of Christ
and guardian of princes sitting high on St.\ Peter's
throne and hurling the thunderbolt of excommunication; no, the beginning was
in meekness, in sacrifice, in patience. The messengers of peace
went out into the world, out to the peoples in their wild
strife, to bring them the Reconciler's greeting and reconcile enemies;
such was the beginning.

What the relation is between power and right can be seen from
this beginning. Had those messengers not been assured
that there is an incorruptible right which is stronger
than brute power, they would not have brought the Reconciler's
\opage{5}message, they would not have ventured out among the wild hordes. They
came defenceless, and yet so well equipped; they came as the
weak, and yet so strong! for he that can rule his own
spirit is strong, stronger than he that can only take the city;
and he that can pray for his enemies is strong, stronger than
he that can only murder his enemies; and he that can sacrifice
his life for truth and right is strong, stronger than he that
can sacrifice nations and kingdoms for imagined honour and power.

So, then, those messengers went, each his own way, out
into the world. A stranger comes to the warriors' camp;
he is unarmed and would contend with the armed, a
weak man who would yet overcome the strong: who in this
strife shall give way? It will not be the unarmed man who gives way
to the armed; no, it will be the armed who
this time must give way to the unarmed man. It is not, after all,
the stranger who trembles at the ruler's glance, grows
afraid of the flashing sword and throws himself down at the king's
feet and begs for protection and mercy; no, he does not tremble
at the ruler's glance; he proclaims his message; he is
not at all afraid of the flashing sword, for he proclaims his
message; he does not beg for mercy, he offers mercy.
Now the wild cries fall silent; it grows still in the camp:
the warrior hears the stranger's message, the king, otherwise
unbending, bows his head, listens and ponders and feels
strangely moved. So it is not the army, this whole
army of strong men, full-blooded warriors, that in this encounter
has the power, the higher power; no, the stranger with his
message, the man of the word with fire in his glance, the pale monk
with the crucifix, it is he who in this encounter has the power,
the higher power; it is he who in this encounter bears
the sign of power, that sign of power which is eternally the sign
of justice. Behold this sign of power! it makes the weak
strong, the strong weak; for the king kneels with his army before
% --- p. 6 ---
\opage{6}the sign of power and of justice. So it is, then, a
judgment, a world-judgment on right and power: the stronger will
have right; the weaker must give way.

There is a difference in right and there is a difference in power;
in this temporal sphere there will always be formed and fabricated, subtly
fabricated, a right that in the eternal sense is nevertheless wrong; in
the lower reality there will always stir a power, a puffed-up, brutal,
obtrusive power, that in the eternal sense is nevertheless impotence.
If the temporal is nonetheless to be compelled to bow
beneath the eternal; if justice itself, the holy thought
that is eternally power, is to take hold of the wheel of the ages and move
reality: then its revelation must be more than a
great vision in a great moment; then its breaking through must
be more than a flare in the cloud, a lightning flash in the night or a
flickering light in the storm; the divine will then not
confine itself to surprising the peoples with signs from heaven,
but will descend to men, condescend to human
purposes, human conditions, and take up its dwelling among
men.

But when the divine thus on its side
stirs, in order to purify, strengthen and transfigure the human,
it cannot fail that something human on its side
will stir as well, in order to weaken, infect and darken the
divine.

With a bold hand, in a popular and practical way, the Catholic Church drew
the supersensible down into the world of the senses; in a popular and
practical way Catholicism let the invisible put on visible
form and gave the holiest of holies bodily, indeed tangible,
presence. It was a step downward; but yet in a human
sense at the same time a step forward. On the sensuous, imaginative
natural man the spirit works through visible, sensuous
means. The people of the Middle Ages, with their ideal dreams, their
deep longings and strong passions, needed spiritual
% --- p. 7 ---
\opage{7}guardians, taskmasters with divine authority;
the age of fist-law, which demanded ordeals, the judgments of God, required to
hear --- God's own voice, required to see --- God's vicar
himself present on earth. It was a peculiar mixture of
barbarism and piety, of defiance and humility, raw sensuality and ascetic self-denial, in countless shades of
error and pain, crimes and remorse and tears,
that in that fabulous age made theocracy necessary.
With theocracy the hierarchy was completed. Thus rose the
high Church; the Church of the Middle Ages on the rubble of antiquity,
on the mighty ruins of heathendom; the Roman Church with its
dim vaults, with its deep chorale and its mystically
solemn language; the papal Church with its all-availing gifts of grace,
an incalculable surplus of good works, an
inestimable treasure of enshrined saints' bones and of all the
merits of the saints. Thus rose the high Church:
it is the spirit of the Middle Ages in the style of the Middle Ages, this
symbolic rite, this cult rich in images and emblems, with
confession and litanies, with altar candles and clouds of incense, priests
and sacraments.

Certainly there was in all this a distortion of the original,
an ever-growing darkening of truth. The light of the
Gospel had long stood in shadow behind the pillars of the Church; the Lord
of the congregation, the Exemplar himself, was forgotten by a host of worshippers
who would rather kneel to the Queen of Heaven. But
despite this darkening of truth the Church's higher
power was still always founded on a higher right: the human
was uplifted by the divine; wondrous powers
stirred in minds and thoughts. The Church in its power and might
had already created for itself a new heaven; there now came a
turning point, when it should truly appear that it was able also to create a new earth; this turning point is the beginning
of the Crusades.

% --- p. 8 ---
\opage{8}That Europe in the eleventh century needed an awakening
is unmistakable. Amid manifold afflictions
and dark forebodings it laboured with heavy, unclear thoughts to
find a way through reality, a goal for deeds,
a possibility of disentangling the confusion of temporal life. It was
a tormented condition. There was indeed a blessedness in heaven
for those who escaped damnation; but on earth there was
nothing good, for the earth was a vale of tears. So
there went a fear through all souls; there were signs in the sun
and the moon; the peoples languished in terror and expectation, for
the powers of heaven were stirring: Doomsday was close at hand.

Should this earthly thing --- once more --- gain real
significance, this could happen only by its coming
directly to serve the heavenly; should there be a
higher goal for endeavours in the temporal, then there had to
be a way on earth, a way of honour, of actions and
of deeds, that led straight to heaven. ``What shall
I do, that I may inherit eternal life?'' It was not a
single rich young man who asked thus in a kind of rashness;
for thousands upon thousands of young men, rich in the strength for deeds
and in knightly courage, each pondered for himself that
question: ``what shall I do, that I may win an
eternal life, an immortal name?'' It was not a little flock of % the Danish does not close this quotation (PRINTED AS IS)
idlers who, because at bottom they wished to do nothing at all,
passed the time with the question: ``what shall we
do, that we might work the work of God?'' it was the peoples of Europe
who, in affliction and want, in ferment and tension,
did not know what they should do.

But the Church knew it; here again the higher power is
a higher right; for the Pope knew what they should do. At
the council in Clermont it was explained to the listening
crowds what Christendom should do. ``Free
Jerusalem!'' so ran the explanation from the Pope's mouth; to free
% --- p. 9 ---
\opage{9}Jerusalem and murder the enemies of the Cross, that was --- for the remission of sins,
for peace of soul and salvation --- the great
knightly deed, the work that God's people were then to do.
``Free Jerusalem!'' so it rang out over the lands of the West; now
the way was found, the goal set, the high goal. ``Free
Jerusalem!'' now it was clear to all that the Pope's thought was
God's thought and the Pope's will God's will; yes, then
the Pope's command must also be that command of power which the hosts
obeyed. ``Free Jerusalem!'' it was an awakening to life,
a resurrection even before the Day of Judgment. % Danish: „Ostandelse“, PRINTED AS IS for Opstandelse

But time passes, experience comes. The first sight
of the lands of the East, of Palestine, of Jerusalem, how
uplifting, beautiful and glorious was this first distant sight in the colours of the imagination
in comparison with all the hard,
gnarled, wretched reality that, step by step, afterwards
forced itself upon the still inexperienced. The more overwrought an
expectation is, the greater is the disappointment that
follows it. The long strife over the grave of the Resurrection is
a romantic epic, a fabulous knightly game with fantastic
alternations of victories and defeats, glittering hopes
and bitter disappointments, flaring enthusiasm and oppressive
dejection. It lasted two hundred years, but then
the age of illusions was over.

But if the age of illusions is over, then it is essentially
over with the Middle Ages, for the Middle Ages are the age of illusions. What is ---
on the way of history --- illusion? The fata morgana of ideas, a mirage of the imagination,
enchanting when in the same vision, in one and the
same image, it can fuse truth with deceit.
So long as the element of truth is still strong enough to master
the deceit, illusion is like a form of revelation for ideal
powers, and in this character can serve enthusiasm precisely;
but as the relation is gradually reversed, as
the deceit gains the upper hand and devours the element of truth,
% --- p. 10 ---
\opage{10}the ideal will withdraw and leave behind its Medusa-like counter-image,
a truly repellent caricature.

Just when the moment had come when the papacy
had exhausted its spiritual powers and so lost its
essential justification, the Pope with the two swords could
indeed turn to South and North and East and West, declare himself
lord of the world, a commanding ruler over peoples and states;
his right, his canon law with glosses and extravagants,
was then brought completely into a system. It is a grand-style
insistence on being in the right that is carried on here; an outrageous strife between power and
right will now begin in earnest.

A merely worldly power must after all always, if it is
in any way to build on a state of law and give itself
even the mere appearance of right, accept certain limits, bind
itself to certain inviolable conditions; but with a spiritual power,
a holy, infallible, divine power, a power that has long
violated the eternal truth and made the holiest of holies a
means to low worldly ends, it is another matter:
here, in the sad sense of the word, the highest wrong becomes the
highest right.

To picture to oneself the Church's spiritual-worldly power in
its extraordinary greatness and extraordinary abuse,
one must consider the almost incalculable mass of usurpations
and arbitrary acts that this power, this alone infallible
power, has been able to impose on nations and princes;
it is a never-interrupted series of scandalous scenes, with which
its holy lust for dominion, its holy avarice, its holy
intrigue and its holy effrontery have managed to fill
the remarkable period that lies between the end of the Crusades
and the beginning of the Reformation.

The stronger will have right; but here there is a difference in
right and a difference in strength. So strong was the Roman Church
in the midst of its corruption that neither the violent storms that
% --- p. 11 ---
\opage{11}raged within it under the papal schism, nor the
many shocks that the humane powers --- the state, science,
the awakening self-esteem of the nations and advancing
culture --- dealt it from without in ever-rising reactions,
were able to subdue the mighty one. One thing, but only
one, was still left that could compel the otherwise invincible; one thing, but only one, was able to break through
that monstrous web of guile, lies and deceit; this
one thing was the truth in all simplicity, the long-derided
truth, the thought of justice, the ideal right that in its
essence is eternally power.

But when a beginning is really to be made here with the ideal
right, that is, originally and from the root, then one cannot begin
here with tangible superiority or with a previously calculated
probability of a decisive victory; if the fight here is to be
for the truth, unconditionally for the eternal right of truth, then
one must here once again begin in weakness, in temporal insignificance,
in outward lowliness: thus the Reformation begins.

``We will have right, for we have power,'' so it sounded in
threats from Rome; ``we shall surely get power, as surely as we
have right,'' so it sounded in reply from Wittenberg. That the holy
right is in its essence eternally power was precisely the thought that animated
Luther; it was a thought of conscience that raised
him high above the fear of men, a thought of God that gave
him strength and courage, such strength, such courage, that
in every moment of decision --- and with how many moments of decision
is his course of life not marked?
--- he always found the decisive word and confirmed the word by decisive
action. At the Diet of Worms a great question of right is
debated; at the Diet of Worms the holders of power are
assembled. Behold the assembly: it is
a numerous, a glittering assembly of knights, prelates,
princes and electors, with the Emperor himself, the newly crowned,
% --- p. 12 ---
\opage{12}in its midst: all that the Holy Roman Empire has to show
of splendour and honour, of pomp and glory, is present here;
and then behold the lowly monk, the harshly accused,
the twice excommunicated, who now enters the assembly
and is to make his defence: what a contrast in the
outward, of lowliness and loftiness, of impotence and power, must
not force itself upon us at this sight; and then --- what a
contrast in the inward, of loftiness and lowliness, of power
and impotence, does not reverse the relation and transform for us the impression
of this sight!

It is, then, once more the outward, tangible power that
threatens to crush the defenceless man, and once more it is the
defenceless man who would contend with the armed, a weak man who
would yet overcome the strong. So Luther would break the yoke
of the lie; so his last answer rang out in the listening assembly:
``Since then Your Imperial Majesty and Your Electoral
and Princely Graces desire a simple, plain and
straightforward answer, I will give one that has neither horns nor
teeth, namely thus: Unless I am overcome and convinced by the testimony
of Holy Scripture or by open, clear and evident
grounds (for I believe neither
the Pope nor the councils, since it is plain as day
that they have often erred and contradicted one another), so that I am persuaded by
the very passages that I have cited and used, and my conscience is captive to the Word of God, I neither can
nor will recant anything, for it is neither safe
nor advisable to do anything against one's conscience. Here
I stand, I can do no other; God help me. Amen!''

It was in the hour of decision a decisive word, and
the word was a decisive act. Who can now doubt
what Luther wants, and what it is for which he fights? For
the holy Scripture of truth, for the open, clear and evident
grounds, and for the incorruptible right of conscience!
% --- p. 13 ---
\opage{13}He stands there, it is true, alone in the assembly of the mighty; but
yet not alone as one who is forsaken by all. Certainly he has
not secured for himself beforehand a physical support, an army of
allied troops that could advance as from an ambush
and come to his rescue at a given signal; but that is no reason
he is without help and aid. In the moment of danger, just
when he acts as one who can no longer negotiate, in
the moment of danger, work is being done in secret for his
cause, in a thousand places, in a thousand ways it is wrought for him,
negotiated for his help and aid. For how many
new thoughts, how many hitherto obscure, suppressed notions, even in crowned heads, what indescribable
presentiments and bold expectations are here not awakened, called
forth as by a stroke of magic by Luther's strong word. And these
new thoughts, these freshly sprung notions, these bold,
enterprising expectations, they stream to meet him, they plead his
cause, they vote for him, they gather in a circle about him,
a mighty, an invisible army, an allied army of clear grounds,
strong sympathies. The lie, it is true, has its sympathies too;
but truth, right and clear grounds --- however
the world may turn --- they will always work for the good
cause; they will gather in the hour of danger, a mighty army
of hidden forces, strong sympathies.

Certainly it was an unequal fight --- of the rays of light
against gloom and darkness, simple truth against a web of
lies; it was unequal fight --- with the weapons of the spirit, the power of the word
alone against physical power and intrigue and cunning
fanaticism; it was unequal fight, but in unequal fight
strength must be tested. And it was tested step by step and
point by point, and the trial was severe; but never
did the holy Scripture of truth fail, and never did
the clear grounds fail, with the strong right. Yes, it was tested,
the trial was heavy; but that sign of justice, that
% --- p. 14 ---
\opage{14}sign which makes the weak strong, the strong weak, was once again
the sign of power on the day of danger, a sign of victory in the storm
of battle, when Luther's voice rang out loud in the storm: ``Hearest thou,
Pope, not thou most holy, no, thou most sinful, that God
from his heaven shall speedily overthrow thy throne and sink it
into the abyss of hell!''

It is not presumption, it is not arrogance, and it is
not blind passions that speak to us through these words of
threat; no, it is Luther's faith, it is his living and rock-firm
conviction that the holy right is in its essence eternally
power. Nor, then, is it the Pope in Rome alone or
the papacy alone at which the threat is aimed; for it is aimed
at every outrageous usurpation; it is aimed at every
arbitrary and despotic power that through a surfeit of lies
and intrigues fancies it can darken everything and defy
everything, both the holy Scripture of truth and the clear grounds and
the original, strong right: upon every power in time so blind, so unjust,
Martin Luther has pronounced a curse.

\begin{center}\rule{3em}{0.4pt}\end{center}

The stronger will have right: it is a judgment, a world-judgment!
But what is a judgment of the world without a Judge of the
world? What then is all our strife about right and power in
this temporal sphere? A waking dream, a play of shadows, if there
is not an eternal right, a holy power, that knows all and
governs all and reveals itself in all as the Judge of the world.

To thy keeping, all-seeing Providence, almighty Judge of the
world, we then commend our beloved fatherland with the father of the land,
our dear King, faithful to his people, King Frederik
the Seventh, and the whole royal house. May thy Spirit, which
is the spirit of wisdom and understanding, the spirit of counsel and might,
be upon the fatherland's faithful men in the King's and in the people's
councils, that the institutions which guard freedom and the
rich, many-sided development of civic life, in spiritual as well as
% --- p. 15 ---
\opage{15}in temporal respects, may be preserved, protected and used, in
an ever wider compass, to reconcile what is at strife and unite what is
divided. And when we further hope and pray that this university
may continue to flourish and never be forsaken by the
living sense for the inquiry into truth and the knowledge of truth,
which alone binds lasting ties between teachers and
hearers: then we think at the same time of every higher
endeavour in church and school, in science and art, that all
the powers of the spirit may at last be directed toward the same goal; we
think of every honest labour, every good and honourable work
here in the land, that we all may have the rich blessing of peace,
so long as it may be granted us to live in peace.

But if peace is impossible, if the hour of decision strikes,
if the call comes to us, when we are to show in deed that we
still --- still at the eleventh hour --- have kept Luther's faith,
his rock-firm conviction that there is an incorruptible right
that in its essence is eternally power: then be thou, Almighty, in
the unequal fight, our shield and buckler and our mighty fortress,
as thou wast Luther's mighty fortress, his shield and buckler!
For it is nothing new that we have before us; ah, no! it is an
old strife, a long lawsuit from time immemorial, from
that Prince Uffo with the rusty sword --- he took up
arms when it counted --- and so on through the changes of the ages, long
ages, down to the last victims who fell, richly
adorned with the roses of battle. It is an old strife,
for out of the mists of antiquity, the twilight of the past, there gathers an % Danish: „Oltids“, PRINTED AS IS for Oldtids
army, a numerous, mighty army, that fights for our right; it
is an allied army of strong memories: O, let it be near
on the day of battle, in the hour of danger; let it break forth with
united strength and with power and fire, like flames in a Luther's
spirit, like the sword's lightning in Gustavus Adolphus's hand, to shine
upon our way to the goal.

\end{document}
